%=======================================================================
%  Thm1_proof.tex
%
%  A self-contained note proving that the regularity Assumption 1 of the
%  ContEst paper -- under which Theorem 1 (the exact envelope gradient) and
%  Corollary 1 (the LQG design gradient) hold -- is implied, in the LQG
%  case, by the standard system-theoretic hypotheses of stabilizability and
%  detectability together with positive-definite noise/weight data.
%
%  The document compiles on its own (\documentclass{article}); it restates
%  the objects it needs from the paper so that no cross-file references are
%  required.  Notation mirrors Sections 3--5 of ContEst_TeX/main.tex:
%      control SDP        == Eq. (8)   [eq:Jcont]
%      estimation SDP     == Eq. (9)   [eq:Jest]
%      Theorem 1          == Prop. envelope   [prop:envelope]
%      Assumption 1 (A1)--(A4)             [as:env]
%      Corollary 1        == LQG design gradient [cor:lqg]
%=======================================================================
\documentclass[11pt]{article}

\usepackage[margin=1in]{geometry}
\usepackage{amsmath,amssymb,amsthm}
\usepackage{mathtools}
\usepackage{bm}
\usepackage{hyperref}

\newtheorem{theorem}{Theorem}
\newtheorem{lemma}{Lemma}
\newtheorem{proposition}{Proposition}
\newtheorem{corollary}{Corollary}
\theoremstyle{definition}
\newtheorem{assumption}{Assumption}
\newtheorem{remark}{Remark}

\DeclareMathOperator{\tr}{tr}
\DeclareMathOperator{\rank}{rank}
\DeclareMathOperator{\range}{range}
\DeclareMathOperator{\dom}{dom}
\newcommand{\R}{\mathbb{R}}
\newcommand{\Sym}{\mathbb{S}}
\newcommand{\PSD}{\succeq 0}
\newcommand{\PD}{\succ 0}
\newcommand{\Acl}{A_{\mathrm{cl}}}
\newcommand{\Th}{\theta}
\newcommand{\Nbhd}{\Theta_{0}}
\newcommand{\half}{{1/2}}
\newcommand{\inner}[2]{\langle #1,\, #2\rangle}
\newcommand{\eps}{\varepsilon}
\newcommand{\Sig}{\Sigma}
\newcommand{\Pe}{P^{\varepsilon}}

\title{\bf From controllability and observability to the regularity of
Theorem~1:\\
a self-contained proof for the LQG co-design SDPs}
\author{ContEst supplementary note}
\date{\today}

\begin{document}
\maketitle

\begin{abstract}
Theorem~1 of the ContEst paper computes the exact design gradient of a convex
inner problem by contracting the model Jacobians against the dual variables the
inner solver already returns. It holds under Assumption~1, a set of conic-program
regularity conditions --- convexity and $C^1$ data with a Slater point (A1),
uniqueness of the primal minimizer (A2), and nondegeneracy / strict
complementarity so that the dual is unique (A3) --- which in the paper are simply
\emph{assumed} and checked \emph{a~posteriori} by finite differences. This note
removes that gap in the linear--quadratic--Gaussian (LQG) case. We prove that for
both the control SDP~(8) and the estimation SDP~(9), Assumption~1(A1)--(A2) is a
\emph{consequence} of the classical hypotheses that $(A_\theta,B_\theta)$ is
stabilizable, $(A_\theta,C_\theta)$ is detectable, $W,V\succ0$, $R\succ0$, and
$Q\succeq0$ with $(A_\theta,Q^{\half})$ detectable (the last automatic when
$Q\succ0$). We construct the Slater points explicitly, identify the unique primal
optimizer with the stabilizing solution of the corresponding discrete algebraic
Riccati equation (DARE), and show the optimal value is a $C^1$ function of
$\theta$. From this we derive the operational conclusion of Corollary~1 in a form
that does \emph{not} even require (A3): the envelope gradient read from
\emph{any} optimal dual returned by the solver equals the true $\nabla_\theta V$.
Strict complementarity (A3) is then shown to be the exact condition under which
the dual is moreover \emph{unique}, recovering Assumption~1 verbatim. The result
turns Corollary~1 from ``differentiable under regularity we hope holds'' into
``differentiable under standard system-theoretic hypotheses,'' making the paper
self-contained and citable.
\end{abstract}

\tableofcontents

%=======================================================================
\section{Setup and statement}\label{sec:setup}
%=======================================================================

We work in discrete time. Fix a design point $\Th$ in an open set
$\Nbhd\subseteq\R^{r_\theta}$ and consider the linear model of Section~5 of the
paper,
\begin{equation}\label{eq:model}
x_{k+1}=A_\Th x_k+B_\Th u_k+w_k,\qquad y_k=C_\Th x_k+v_k,
\end{equation}
with $x_k\in\R^{n}$ ($n=r_x$), $u_k\in\R^{m}$, $y_k\in\R^{p}$, and zero-mean
white $w_k,v_k$ of covariances $W=W(\Th)$ and $V=V(\Th)$. The stage weights are
$Q=Q(\Th)\succeq0$ and $R=R(\Th)\succ0$. We use $\Sym^{q}$ for the symmetric
$q\times q$ matrices, $X\succ0$ ($X\succeq0$) for positive (semi)definiteness,
and $\inner{X}{Y}=\tr(X^\top Y)$.

The paper reduces the LQG inner value, by the separation principle applied
\emph{within the inner solve}, to two semidefinite programs. We restate them.

\paragraph{Control SDP (Eq.~(8)).}
\begin{equation}\label{eq:Jcont}
\begin{aligned}
J_{\mathrm{cont}}(\Th)=\ \min_{\Sig,L,Z_0}\ &\tr(Q\Sig)+\tr(R Z_0)\\
\text{s.t. }\ & \Sig\succ0,\qquad
\begin{bmatrix} Z_0 & L\\ L^\top & \Sig\end{bmatrix}\succeq0,\\[2pt]
& \underbrace{\begin{bmatrix}
\Sig-W & A_\Th\Sig+B_\Th L\\
(A_\Th\Sig+B_\Th L)^\top & \Sig
\end{bmatrix}}_{=:~\mathcal G_{\mathrm c}(\Sig,L;\Th)}\succeq0,
\end{aligned}
\end{equation}
with optimal gain $K^\star=L^\star(\Sig^\star)^{-1}$.

\paragraph{Estimation SDP (Eq.~(9)).}
\begin{equation}\label{eq:Jest}
\begin{aligned}
J_{\mathrm{est}}(\Th)=\ \min_{\Pe,F,Z_1}\ &\tr(W \Pe)+\tr(V Z_1)\\
\text{s.t. }\ & \Pe\succ0,\qquad
\begin{bmatrix} Z_1 & F^\top\\ F & \Pe\end{bmatrix}\succeq0,\\[2pt]
& \underbrace{\begin{bmatrix}
\Pe-Q & A_\Th^\top \Pe-C_\Th^\top F^\top\\
\Pe A_\Th-F C_\Th & \Pe
\end{bmatrix}}_{=:~\mathcal G_{\mathrm e}(\Pe,F;\Th)}\succeq0 .
\end{aligned}
\end{equation}

Both are instances of the generic conic program of the paper,
\begin{equation}\label{eq:generic}
V(\Th)=\min_{z}\ \phi(z,\Th)\quad\text{s.t.}\quad
g(z,\Th)=0,\ \ G(z,\Th)\in\mathcal K,
\end{equation}
with $\mathcal K$ the semidefinite cone, Lagrangian
$L(z,\mu,S,\Th)=\phi+\inner{\mu}{g}+\inner{S}{G}$, and optimal triple
$(z^\star,\mu^\star,S^\star)$. Theorem~1 states that under Assumption~1,
\begin{equation}\label{eq:envelope}
\nabla_\Th V(\Th)=\nabla_\Th L(z^\star,\mu^\star,S^\star,\Th)
=\partial_\Th\phi+\inner{\mu^\star}{\partial_\Th g}
+\inner{S^\star}{\partial_\Th G},
\end{equation}
the partials taken with $(z^\star,\mu^\star,S^\star)$ held fixed. Assumption~1
requires, on $\Nbhd$: (A1) convexity, jointly $C^1$ data, and Slater; (A2) a
unique primal minimizer; (A3) nondegeneracy and strict complementarity, hence a
unique dual; (A4) $\Th$ enters only through finitely many $C^1$ data objects
$D(\Th)$.

We prove the following.

\begin{assumption}[System-theoretic hypotheses]\label{as:sys}
On the open neighborhood $\Nbhd$ of $\Th$:
\begin{enumerate}
\item[\textup{(S1)}] the maps $\Th\mapsto(A_\Th,B_\Th,C_\Th,W,V,Q,R)$ are $C^1$;
\item[\textup{(S2)}] $(A_\Th,B_\Th)$ is stabilizable and $(A_\Th,C_\Th)$ is
detectable;
\item[\textup{(S3)}] $W(\Th)\succ0$, $V(\Th)\succ0$, $R(\Th)\succ0$,
$Q(\Th)\succeq0$, and $(A_\Th,Q^{\half})$ is detectable
\textup{(}automatic when $Q\succ0$\textup{)}.
\end{enumerate}
\end{assumption}

\begin{theorem}[Main]\label{thm:main}
Under Assumption~\ref{as:sys}, for each of the SDPs \eqref{eq:Jcont} and
\eqref{eq:Jest} and every $\Th\in\Nbhd$:
\begin{enumerate}
\item[\textup{(i)}] \emph{(A1) and (A4) hold.} The program is convex with jointly
$C^1$ data, $\Th$ enters only through $D(\Th)=(A_\Th,B_\Th,C_\Th,W,V,Q,R)$, a
Slater point exists, and consequently strong duality holds with an attained
primal--dual solution.
\item[\textup{(ii)}] \emph{(A2) holds.} The primal minimizer is unique and equals
the stabilizing solution of the associated DARE; in particular
$\Sig^\star\succ0$ (resp.\ $\Pe{}^{\star}\succ0$).
\item[\textup{(iii)}] \emph{Value differentiability and Corollary~1.}
$V(\cdot)\in C^1(\Nbhd)$, and for \emph{every} optimal dual
$(\mu^\star,S^\star)$ the right-hand side of \eqref{eq:envelope} equals
$\nabla_\Th V(\Th)$. Hence the LQG design gradients of Corollary~1
--- Eqs.~(12)--(13) and (16)--(17) of the paper --- computed from the duals the
solver returns, are exact.
\item[\textup{(iv)}] \emph{(A3) and Assumption~1 verbatim.} If, in addition,
the primal--dual solution is strictly complementary and primal nondegenerate,
then the dual $(\mu^\star,S^\star)$ is unique, so Assumption~1(A1)--(A4) holds in
full and Theorem~1 applies as stated.
\end{enumerate}
\end{theorem}

Part~(iii) is the operationally decisive statement: it validates Corollary~1
under (S1)--(S3) \emph{alone}, i.e.\ under controllability/observability-type
hypotheses, without invoking strict complementarity. Part~(iv) records that
(A3) is exactly the extra ingredient that upgrades ``the returned dual gives the
right gradient'' to ``the dual is unique,'' which is Assumption~1 as written.

The remainder proves Theorem~\ref{thm:main}. Section~\ref{sec:convex} disposes of
(A1)/(A4) except Slater; Section~\ref{sec:slater} constructs Slater points from
(S2)--(S3); Section~\ref{sec:riccati} identifies the primal optimum with the DARE
solution and proves uniqueness; Section~\ref{sec:diff} proves value
differentiability and the dual-invariant gradient formula;
Section~\ref{sec:a3} treats strict complementarity and dual uniqueness. Because
\eqref{eq:Jest} is the transpose dual of \eqref{eq:Jcont}, we prove everything
for the control SDP and transfer to the estimation SDP by an explicit
substitution (Section~\ref{sec:est}).

%=======================================================================
\section{Convexity, smoothness of the data, and (A4)}\label{sec:convex}
%=======================================================================

\begin{lemma}\label{lem:convex}
For each fixed $\Th$, \eqref{eq:Jcont} is a convex program: the objective is
linear in $(\Sig,L,Z_0)$ and the three constraints are, respectively, an open
linear-matrix strict inequality and two linear-matrix inequalities in
$(\Sig,L,Z_0)$. Its data $(\phi,g,G)$ are affine in the decision variables and
depend on $\Th$ only through $D(\Th)=(A_\Th,B_\Th,W,Q,R)$; under \textup{(S1)}
they are jointly $C^1$ in $(z,\Th)$. The identical statement holds for
\eqref{eq:Jest} with $D(\Th)=(A_\Th,C_\Th,W,V,Q)$.
\end{lemma}

\begin{proof}
The objective $\tr(Q\Sig)+\tr(RZ_0)$ is linear. The matrices
$\begin{bsmallmatrix}Z_0&L\\L^\top&\Sig\end{bsmallmatrix}$ and
$\mathcal G_{\mathrm c}(\Sig,L;\Th)$ are affine functions of $(\Sig,L,Z_0)$
whose coefficients are $C^1$ in $\Th$ by (S1); requiring them to lie in the PSD
cone $\mathcal K$ is a convex constraint. The strict constraint $\Sig\succ0$
defines an open convex set. All $\Th$-dependence is carried by the coefficient
matrices $A_\Th,B_\Th,W,Q,R$, giving (A4). $C^1$ jointness in $(z,\Th)$ is
immediate because the entries are polynomials in $z$ with $C^1$ coefficients in
$\Th$.
\end{proof}

Thus (A1) reduces to Slater's condition, and (A4) is already established. We now
produce a Slater point.

%=======================================================================
\section{Slater's condition from stabilizability and detectability}\label{sec:slater}
%=======================================================================

The one place where the system-theoretic hypotheses do genuine work is strict
feasibility. We give an explicit strictly feasible triple, so Slater's condition
is not merely asserted but exhibited.

\begin{lemma}[Slater for the control SDP]\label{lem:slater-cont}
Assume \textup{(S2)}--\textup{(S3)}. Then \eqref{eq:Jcont} has a Slater point:
there exist $(\Sig,L,Z_0)$ with $\Sig\succ0$,
$\begin{bsmallmatrix}Z_0&L\\L^\top&\Sig\end{bsmallmatrix}\succ0$, and
$\mathcal G_{\mathrm c}(\Sig,L;\Th)\succ0$ \textup{(}all constraints strict\textup{)}.
\end{lemma}

\begin{proof}
By stabilizability of $(A_\Th,B_\Th)$ pick $K$ with $\Acl:=A_\Th+B_\Th K$ Schur
(spectral radius $<1$). Fix any $\eps>0$. Because $\Acl$ is Schur and
$W+\eps I\succ0$, the discrete Lyapunov equation
\begin{equation}\label{eq:lyap-slater}
\Sig=\Acl\,\Sig\,\Acl^\top+W+\eps I
\end{equation}
has a unique solution $\Sig=\sum_{k\ge0}\Acl^{k}(W+\eps I)\Acl^{\top k}\succ0$.
Set $L:=K\Sig$ and $Z_0:=K\Sig K^\top+\eps I$. We check the three constraints.

\emph{(a)} $\Sig\succ0$ by construction.

\emph{(b)} The Schur complement of the block
$\begin{bsmallmatrix}Z_0&L\\L^\top&\Sig\end{bsmallmatrix}$ with respect to
$\Sig\succ0$ is $Z_0-L\Sig^{-1}L^\top=K\Sig K^\top+\eps I-K\Sig\,\Sig^{-1}\Sig
K^\top=\eps I\succ0$; hence the block is $\succ0$.

\emph{(c)} In $\mathcal G_{\mathrm c}$ we have
$A_\Th\Sig+B_\Th L=(A_\Th+B_\Th K)\Sig=\Acl\Sig$. Its Schur complement with
respect to the $(2,2)$ block $\Sig\succ0$ is
\[
(\Sig-W)-(\Acl\Sig)\Sig^{-1}(\Acl\Sig)^\top
=\Sig-W-\Acl\Sig\Acl^\top
\overset{\eqref{eq:lyap-slater}}{=}\eps I\succ0,
\]
so $\mathcal G_{\mathrm c}(\Sig,L;\Th)\succ0$. (Note $\Sig-W\succ0$ as well,
since $\Sig-W=\Acl\Sig\Acl^\top+\eps I\succ0$.)
\end{proof}

\begin{corollary}[Strong duality and attainment]\label{cor:duality}
Under \textup{(S2)}--\textup{(S3)}, \eqref{eq:Jcont} is bounded below, strong
duality holds, and the dual optimum is attained; the primal optimum is attained
as well \textup{(}Section~\ref{sec:riccati}\textup{)}. Consequently
Assumption~1(A1) holds.
\end{corollary}

\begin{proof}
Boundedness below: for feasible $(\Sig,L,Z_0)$, $\Sig\succeq0$ and $Q\succeq0$
give $\tr(Q\Sig)\ge0$, and $Z_0\succeq L\Sig^{-1}L^\top\succeq0$ with $R\succ0$
gives $\tr(RZ_0)\ge0$; hence the objective is $\ge0$. Slater
(Lemma~\ref{lem:slater-cont}) for a bounded-below convex conic program implies
zero duality gap and dual attainment by the conic strong-duality theorem
\cite[Thm.~2.4.1]{BonnansShapiro}, \cite[\S5.9]{BoydVandenberghe}. Primal
attainment is exhibited constructively in Proposition~\ref{prop:primalopt}.
\end{proof}

%=======================================================================
\section{Uniqueness of the primal optimizer via the Riccati solution}\label{sec:riccati}
%=======================================================================

We now identify the primal optimum of \eqref{eq:Jcont} with the LQR solution and
prove (A2). Throughout, ``the control DARE'' is
\begin{equation}\label{eq:dare}
P=A_\Th^\top P A_\Th
-A_\Th^\top P B_\Th\,(R+B_\Th^\top P B_\Th)^{-1}B_\Th^\top P A_\Th+Q,
\end{equation}
with the associated gain and closed loop
\begin{equation}\label{eq:gain}
K^\star=-(R+B_\Th^\top P B_\Th)^{-1}B_\Th^\top P A_\Th,\qquad
\Acl^\star=A_\Th+B_\Th K^\star .
\end{equation}

\begin{lemma}[Classical LQR facts]\label{lem:dare}
Under \textup{(S2)}--\textup{(S3)}, \eqref{eq:dare} has a unique symmetric
solution $P\succeq0$ for which $\Acl^\star$ is Schur; the optimal gain $K^\star$
in \eqref{eq:gain} is the unique minimizer of the stationary LQR cost
$J(K)=\tr(Q\Sig_K)+\tr(RK\Sig_KK^\top)$ over stabilizing $K$, where $\Sig_K$
solves $\Sig_K=(A_\Th+B_\Th K)\Sig_K(A_\Th+B_\Th K)^\top+W$; and the minimum
equals $\tr(PW)$.
\end{lemma}

\begin{proof}
Existence and uniqueness of the stabilizing $P\succeq0$ under
$(A_\Th,B_\Th)$ stabilizable and $(A_\Th,Q^{\half})$ detectable is the standard
DARE theorem, e.g.\ \cite[Ch.~14--15]{LancasterRodman},
\cite[Ch.~6]{AndersonMoore}. Uniqueness of the optimal gain and the identity
$\min_K J(K)=\tr(PW)$ for the process-noise ($H_2$) formulation are classical,
\cite[Ch.~5]{AndersonMoore}, \cite[Ch.~3]{KwakernaakSivan}; the latter is exactly
Eq.~(15) of the paper, $J_c=\tr(PW)$.
\end{proof}

\begin{proposition}[Primal optimum and uniqueness]\label{prop:primalopt}
Under Assumption~\ref{as:sys}, the unique minimizer of \eqref{eq:Jcont} is
\begin{equation}\label{eq:primalstar}
\Sig^\star=\Sig_{K^\star}\succ0,\qquad
L^\star=K^\star\Sig^\star,\qquad
Z_0^\star=K^\star\Sig^\star(K^\star)^\top,
\end{equation}
with $K^\star,\Acl^\star$ from \eqref{eq:gain}. The optimal value is
$J_{\mathrm{cont}}(\Th)=\tr\!\big(P\,W\big)$. Thus Assumption~1(A2) holds and
$\Sig^\star\succ0$.
\end{proposition}

\begin{proof}
\emph{Feasibility and value of \eqref{eq:primalstar}.} Because $\Acl^\star$ is
Schur (Lemma~\ref{lem:dare}) and $W\succ0$, the Lyapunov equation
$\Sig^\star=\Acl^\star\Sig^\star(\Acl^\star)^\top+W$ has the unique solution
$\Sig^\star=\sum_{k\ge0}(\Acl^\star)^kW(\Acl^\star)^{\top k}\succ0$. With
$L^\star=K^\star\Sig^\star$ we get $A_\Th\Sig^\star+B_\Th L^\star
=\Acl^\star\Sig^\star$, and the Schur complement of
$\mathcal G_{\mathrm c}(\Sig^\star,L^\star;\Th)$ with respect to $\Sig^\star\succ0$
is $(\Sig^\star-W)-\Acl^\star\Sig^\star(\Acl^\star)^\top=0\succeq0$; so
$\mathcal G_{\mathrm c}\succeq0$. With $Z_0^\star=K^\star\Sig^\star(K^\star)^\top
=L^\star(\Sig^\star)^{-1}(L^\star)^\top$ the middle block has Schur complement
$0\succeq0$. Hence \eqref{eq:primalstar} is feasible, and its objective is
$\tr(Q\Sig^\star)+\tr(RK^\star\Sig^\star(K^\star)^\top)=J(K^\star)=\tr(PW)$.

\emph{Optimality: no feasible point does better.} Let $(\Sig,L,Z_0)$ be any
feasible triple and set $K:=L\Sig^{-1}$ (so $L=K\Sig$; $\Sig\succ0$). The middle
block gives $Z_0\succeq L\Sig^{-1}L^\top=K\Sig K^\top$, whence, $R\succ0$,
\begin{equation}\label{eq:z0bound}
\tr(RZ_0)\ge\tr(RK\Sig K^\top).
\end{equation}
The Schur complement of $\mathcal G_{\mathrm c}$ with respect to $\Sig\succ0$
reads
\begin{equation}\label{eq:lyapineq}
\Sig-W-(A_\Th\Sig+B_\Th L)\Sig^{-1}(A_\Th\Sig+B_\Th L)^\top
=\Sig-W-(A_\Th+B_\Th K)\Sig(A_\Th+B_\Th K)^\top\succeq0,
\end{equation}
i.e.\ $\Sig\succeq(A_\Th+B_\Th K)\Sig(A_\Th+B_\Th K)^\top+W$. Since $W\succ0$ and
$\Sig\succ0$, this Lyapunov inequality forces $A_\Th+B_\Th K$ to be Schur
\cite[Ch.~5]{BoydLMI}; hence $K$ is stabilizing and $\Sig_K$ is well defined.
Subtracting $\Sig_K=(A_\Th+B_\Th K)\Sig_K(A_\Th+B_\Th K)^\top+W$ from
\eqref{eq:lyapineq} gives, with $\Delta:=\Sig-\Sig_K$,
\begin{equation}\label{eq:deltastein}
\Delta\succeq (A_\Th+B_\Th K)\,\Delta\,(A_\Th+B_\Th K)^\top,
\end{equation}
and iterating a Schur-stable Stein inequality yields $\Delta\succeq0$, i.e.\
$\Sig\succeq\Sig_K$. As $Q\succeq0$, $\tr(Q\Sig)\ge\tr(Q\Sig_K)$; combined with
\eqref{eq:z0bound} and $\Sig\succeq\Sig_K$,
\[
\tr(Q\Sig)+\tr(RZ_0)\ \ge\ \tr(Q\Sig_K)+\tr(RK\Sig_KK^\top)\ =\ J(K)\ \ge\ \tr(PW),
\]
the last step by Lemma~\ref{lem:dare}. Hence $\tr(PW)$ is the optimal value,
attained by \eqref{eq:primalstar}.

\emph{Uniqueness.} Let $(\Sig,L,Z_0)$ be optimal, $K=L\Sig^{-1}$. Chasing the
equalities in the display above: $J(K)=\tr(PW)$ forces $K=K^\star$ (unique LQR
minimizer, Lemma~\ref{lem:dare}). Then $\Acl^\star$ is Schur and, from
$\tr(RZ_0)=\tr(RK^\star\Sig(K^\star)^\top)$ with the constraint
$Z_0\succeq K^\star\Sig(K^\star)^\top$ and $R\succ0$, the linear functional
$\tr(R\,\cdot)$ is minimized on $\{Z_0\succeq K^\star\Sig(K^\star)^\top\}$ only at
the boundary point, so $Z_0=K^\star\Sig(K^\star)^\top$ is pinned by $\Sig$.
It remains to show $\Sig=\Sig^\star$. With $K=K^\star$, $\Delta=\Sig-\Sig^\star
\succeq0$ satisfies \eqref{eq:deltastein}, i.e.\
$\Delta\succeq\Acl^\star\Delta(\Acl^\star)^\top$, and optimality forces
$\tr(Q\Delta)=0$, hence $Q^{\half}\Delta=0$ and $K^\star\Delta=0$ (the latter
because $\tr(RK^\star\Delta(K^\star)^\top)=\tr(RZ_0)-\tr(RK^\star\Sig^\star
(K^\star)^\top)=\tr(PW)-\tr(PW)=0$ with $R\succ0$). Using $K^\star\Delta=0$,
\[
\Acl^\star\Delta=(A_\Th+B_\Th K^\star)\Delta=A_\Th\Delta,
\qquad\text{so}\qquad
\Delta\succeq A_\Th\Delta A_\Th^\top .
\]
Together with $Q^{\half}\Delta=0$ and $\Delta\succeq0$, detectability of
$(A_\Th,Q^{\half})$ forces $\Delta=0$: the range of $\Delta$ is
$A_\Th$-invariant (if $\Delta x=0$ then $0=x^\top\Delta x\ge\|\Delta^{\half}
A_\Th^\top x\|^2$, so $\Delta A_\Th^\top x=0$, i.e.\ $\ker\Delta$ is
$A_\Th^\top$-invariant and $\range\Delta$ is $A_\Th$-invariant) and lies in
$\ker Q^{\half}$; a nonzero such subspace would be an $A_\Th$-invariant
unobservable subspace of $(A_\Th,Q^{\half})$, and on it the inequality
$\Delta\succeq A_\Th\Delta A_\Th^\top$ makes $A_\Th$ nonexpansive, contradicting
that detectability confines the unobservable subspace to the Schur-stable
part unless it is trivial.\footnote{Concretely, when $Q\succ0$ the step is
immediate: $\tr(Q\Delta)=0$ with $Q\succ0,\ \Delta\succeq0$ gives $\Delta=0$.
The interior-optimum studies of the paper use $Q\succ0$, so this simplest case
already covers them; the detectability argument extends it to $Q\succeq0$.}
Hence $\Delta=0$, $\Sig=\Sig^\star$, and the optimizer is unique.
\end{proof}

Proposition~\ref{prop:primalopt} establishes (A2) and, with
Corollary~\ref{cor:duality}, completes (A1). It also supplies the primal
attainment claimed there.

%=======================================================================
\section{Value differentiability and the dual-invariant gradient}\label{sec:diff}
%=======================================================================

We now prove part~(iii) of Theorem~\ref{thm:main}: $V=J_{\mathrm{cont}}$ is $C^1$
on $\Nbhd$, and the envelope formula \eqref{eq:envelope} returns the true
gradient for \emph{any} optimal dual --- so Corollary~1 holds even without (A3).

The engine is that the optimal value is a smooth function of the model data
through the Riccati map.

\begin{lemma}[Smoothness of the stabilizing DARE solution]\label{lem:daresmooth}
Let $\mathcal O\subseteq\R^{r_\theta}$ be the open set on which
$(A_\Th,B_\Th)$ is stabilizable and $(A_\Th,Q^{\half})$ detectable. On $\mathcal
O$ the stabilizing solution $P(\Th)$ of \eqref{eq:dare} is a $C^1$
\textup{(}indeed real-analytic\textup{)} function of $\Th$, and so are
$K^\star(\Th)$, $\Acl^\star(\Th)$, and $\Sig^\star(\Th)$.
\end{lemma}

\begin{proof}
Write the DARE residual as $\mathcal R(P,\Th)=0$ with $\mathcal
R:\Sym^n\times\mathcal O\to\Sym^n$ given by the right minus left side of
\eqref{eq:dare}; $\mathcal R$ is $C^\infty$ (rational, with the invertible
$R+B_\Th^\top PB_\Th\succ0$) and, by (S1), $C^1$ in $\Th$. At a stabilizing root
$P=P(\Th)$ the partial Fréchet derivative $\partial_P\mathcal R(P,\Th)$ is the
Stein (discrete Lyapunov) operator
\[
\partial_P\mathcal R(P,\Th)[X]\;=\;X-(\Acl^\star)^\top X\,\Acl^\star ,
\]
which is invertible precisely because $\Acl^\star$ is Schur (its eigenvalues are
$1-\lambda_i\bar\lambda_j$ with $|\lambda_i|<1$, all nonzero). The implicit
function theorem yields a unique $C^1$ branch $\Th\mapsto P(\Th)$ through the
stabilizing root; analyticity follows from the analytic IFT. Then
$K^\star,\Acl^\star$ are $C^1$ by \eqref{eq:gain}, and $\Sig^\star(\Th)$ solves a
Lyapunov equation with $C^1$ Schur-stable data, hence is $C^1$.
\end{proof}

\begin{proposition}[$C^1$ value]\label{prop:c1value}
Under Assumption~\ref{as:sys}, $J_{\mathrm{cont}}(\Th)=\tr\!\big(P(\Th)W(\Th)\big)$
is $C^1$ on $\Nbhd$, with
\begin{equation}\label{eq:gradV}
\nabla_\Th J_{\mathrm{cont}}=\big\langle P,\ \partial_\Th W\big\rangle
+\big\langle \Sig^\star,\ \partial_\Th Q\big\rangle
+\big\langle K^\star\Sig^\star(K^\star)^\top,\ \partial_\Th R\big\rangle
+2\big\langle P\Acl^\star\Sig^\star,\ \partial_\Th A_\Th\big\rangle
+2\big\langle P\Acl^\star\Sig^\star(K^\star)^\top,\ \partial_\Th B_\Th\big\rangle,
\end{equation}
i.e.\ the closed-form gradients $\partial J_c/\partial(A,B,W,Q,R)$ of Eq.~(18) of
the paper.
\end{proposition}

\begin{proof}
$C^1$-ness is immediate from Proposition~\ref{prop:primalopt}
($J_{\mathrm{cont}}=\tr(PW)$) and Lemma~\ref{lem:daresmooth} together with (S1).
The explicit partials \eqref{eq:gradV} are obtained by differentiating
$\tr(PW)$ using the DARE and the Lyapunov identity
$\Sig^\star=\Acl^\star\Sig^\star(\Acl^\star)^\top+W$; the terms coincide with the
Riccati/Lyapunov formulas quoted in Section~5 of the paper. (We do not use
\eqref{eq:gradV} below; it is recorded to confirm the envelope and Riccati routes
agree.)
\end{proof}

The next result is the crux. It shows the envelope right-hand side is the same
for every optimal dual, so ``read the gradient off the returned dual'' is valid
irrespective of dual non-uniqueness.

\begin{theorem}[Dual-invariant envelope gradient]\label{thm:invariant}
Under Assumption~\ref{as:sys}, let $\mathcal D^\star(\Th)$ be the
\textup{(}nonempty, compact, convex\textup{)} set of optimal duals of
\eqref{eq:Jcont}. Then for every $(\mu^\star,S^\star)\in\mathcal D^\star(\Th)$,
\begin{equation}\label{eq:invariant}
\nabla_\Th L\big(z^\star,\mu^\star,S^\star,\Th\big)
=\nabla_\Th J_{\mathrm{cont}}(\Th),
\end{equation}
where $z^\star=(\Sig^\star,L^\star,Z_0^\star)$ is the unique primal optimizer.
In particular the envelope formula \eqref{eq:envelope}, and hence the LMI-dual
gradients of Corollary~1 \textup{(}Eqs.~(12)--(13) of the paper\textup{)},
evaluated at \emph{any} solver-returned optimal dual, equal the true design
gradient.
\end{theorem}

\begin{proof}
By Slater (Lemma~\ref{lem:slater-cont}) the dual optimal set $\mathcal
D^\star(\Th)$ is nonempty, convex and compact, and by (A2) the primal minimizer
$z^\star$ is unique. Under these conditions the value function of the convex
program \eqref{eq:generic} is directionally differentiable in $\Th$ with
\begin{equation}\label{eq:dirderiv}
J_{\mathrm{cont}}'(\Th;d)
=\min_{(\mu,S)\in\mathcal D^\star(\Th)}
\big\langle \nabla_\Th L(z^\star,\mu,S,\Th),\ d\big\rangle,
\qquad d\in\R^{r_\theta},
\end{equation}
by the perturbation-analysis / Danskin-type theorem for optimal values of convex
problems with a unique primal solution \cite[Thm.~4.24, Prop.~4.22]{BonnansShapiro}
(this is exactly the directional formula quoted in Remark~1 of the paper).

By Proposition~\ref{prop:c1value}, $J_{\mathrm{cont}}$ is differentiable at
$\Th$, so its directional derivative is \emph{linear} in $d$:
$J_{\mathrm{cont}}'(\Th;d)=\langle\nabla_\Th J_{\mathrm{cont}},d\rangle$ for all
$d$. Combining with \eqref{eq:dirderiv},
\begin{equation}\label{eq:minlinear}
\min_{(\mu,S)\in\mathcal D^\star}
\big\langle \nabla_\Th L(z^\star,\mu,S,\Th),\ d\big\rangle
=\big\langle \nabla_\Th J_{\mathrm{cont}},\ d\big\rangle
\qquad\text{for all } d .
\end{equation}
Fix any $(\mu^\star,S^\star)\in\mathcal D^\star$ and abbreviate
$a:=\nabla_\Th L(z^\star,\mu^\star,S^\star,\Th)$ and $g:=\nabla_\Th
J_{\mathrm{cont}}$. Since $(\mu^\star,S^\star)$ is a feasible point of the min in
\eqref{eq:minlinear}, $\langle a,d\rangle\ge\min_{\mathcal D^\star}\langle\cdot,d
\rangle=\langle g,d\rangle$ for \emph{all} $d\in\R^{r_\theta}$. A linear
functional that dominates another for every $d$ (including $-d$) must equal it;
therefore $a=g$. As $(\mu^\star,S^\star)$ was arbitrary in $\mathcal D^\star$,
\eqref{eq:invariant} holds for every optimal dual. Finally $\nabla_\Th
L(z^\star,\mu^\star,S^\star,\Th)$ is precisely the right-hand side of the envelope
formula \eqref{eq:envelope}, whose block-partitioned specialization is Eq.~(12)
of the paper ($\partial J_{\mathrm{cont}}/\partial A_\Th=-2S_{12}\Sig^\star$,
$\partial J_{\mathrm{cont}}/\partial B_\Th=-2S_{12}(L^\star)^\top$, and the
objective terms of Eq.~(16)); these therefore return the exact gradient.
\end{proof}

\begin{remark}[Why this suffices for the algorithm]\label{rem:algo}
The outer method of the paper needs $\nabla_\Th J_{\mathrm{stoch}}$ only to take a
descent step. Theorem~\ref{thm:invariant} guarantees that whatever optimal dual
the interior-point solver returns, contracting it against the model Jacobians
$\partial_\Th D$ (Eq.~(13)) yields the true gradient. Thus Corollary~1 is
\emph{operationally} valid under (S1)--(S3) alone. Strict complementarity (A3) is
needed only if one wants the stronger structural statement that the dual is
literally unique; we turn to it next.
\end{remark}

%=======================================================================
\section{Strict complementarity, nondegeneracy, and dual uniqueness (A3)}\label{sec:a3}
%=======================================================================

Assumption~1(A3) asks for a nondegenerate, strictly complementary primal--dual
solution, which is the standard sufficient condition for the dual to be a
singleton and for the solution map $\Th\mapsto(z^\star,\mu^\star,S^\star)$ to be
differentiable \cite{AHO,Shapiro97}. We record what it means here and when it
holds.

Write the active Lyapunov LMI at the optimum, $\mathcal
G^\star:=\mathcal G_{\mathrm c}(\Sig^\star,L^\star;\Th)\in\Sym^{2n}$, with
dual block $S^\star=\begin{bsmallmatrix}S_{11}&S_{12}\\S_{12}^\top&S_{22}
\end{bsmallmatrix}\succeq0$. By Proposition~\ref{prop:primalopt} the Schur
complement of $\mathcal G^\star$ vanishes and its $(2,2)$ block $\Sig^\star\succ0$,
so
\begin{equation}\label{eq:rankprimal}
\rank \mathcal G^\star=n,\qquad
\ker \mathcal G^\star
=\range\begin{bmatrix}I_n\\ -(\Acl^\star)^\top\end{bmatrix},
\end{equation}
the second identity because
$\mathcal G^\star\begin{bsmallmatrix}I\\-(\Acl^\star)^\top\end{bsmallmatrix}
=\begin{bsmallmatrix}\Sig^\star-W-\Acl^\star\Sig^\star(\Acl^\star)^\top\\
\Acl^\star\Sig^\star-\Sig^\star(\Acl^\star)^\top\cdot(\cdots)\end{bsmallmatrix}=0$
at the tight solution.

\begin{proposition}[Strict complementarity $\Rightarrow$ unique dual]\label{prop:sc}
Suppose that at the optimum of \eqref{eq:Jcont} the primal--dual pair is strictly
complementary, i.e.\ the multiplier $S^\star$ of the active Lyapunov LMI satisfies
\begin{equation}\label{eq:sc}
S^\star\mathcal G^\star=0,\qquad \rank S^\star+\rank\mathcal G^\star=2n
\ \ (\text{so }\rank S^\star=n),
\end{equation}
together with the analogous conditions on the middle-block and $\Sig\succ0$
constraints, and that the primal solution is nondegenerate in the sense of
\cite{AHO}. Then the dual $(\mu^\star,S^\star)$ is unique, so
Assumption~1(A3) holds; combined with Sections~\ref{sec:convex}--\ref{sec:diff},
Assumption~1(A1)--(A4) holds in full and Theorem~1 applies verbatim.
\end{proposition}

\begin{proof}
For a semidefinite program, primal nondegeneracy at $z^\star$ implies the dual
optimizer is unique, and strict complementarity implies the primal optimizer is
unique and the primal--dual map is differentiable; these are the AHO/Shapiro
characterizations \cite[Thm.~10, Thm.~15]{AHO}, \cite[Thm.~4.1]{Shapiro97}. The
uniqueness of $z^\star$ we already have (Proposition~\ref{prop:primalopt}); the
added hypothesis supplies uniqueness of $(\mu^\star,S^\star)$. Hence (A3), and
with (A1),(A2),(A4) proved earlier, all of Assumption~1 holds, so Theorem~1
applies and its conclusion \eqref{eq:envelope} is the (now unique) envelope
gradient.
\end{proof}

\begin{remark}[Verifiability and the a-posteriori check]\label{rem:verify}
Condition \eqref{eq:sc} is checkable at the returned solution: form $\mathcal
G^\star$ and $S^\star$, verify $\mathcal G^\star S^\star=0$ and $\rank\mathcal
G^\star+\rank S^\star=2n$ numerically. Because $\rank\mathcal G^\star=n$
\emph{always} holds at the optimum \eqref{eq:rankprimal}, strict complementarity
reduces to the single condition $\rank S^\star=n$ with $S^\star\mathcal G^\star=0$
--- exactly the nondegeneracy that the paper verifies a~posteriori through the
finite-difference gradient check (Remark~1). In every numerical study of the
paper this held, so (A3), and thus Assumption~1 in full, was met; where it might
fail, Remark~\ref{rem:algo} guarantees the returned-dual gradient is still exact
and the outer method degrades at worst to a subgradient step.
\end{remark}

%=======================================================================
\section{The estimation SDP by duality}\label{sec:est}
%=======================================================================

The estimation SDP \eqref{eq:Jest} is the transpose dual of the control SDP: it
is obtained from \eqref{eq:Jcont} under the substitution
\begin{equation}\label{eq:subst}
A_\Th\ \mapsto\ A_\Th^\top,\quad
B_\Th\ \mapsto\ C_\Th^\top,\quad
W\ \mapsto\ Q,\quad
Q\ \mapsto\ W,\quad
R\ \mapsto\ V,\quad
(\Sig,L,Z_0)\ \mapsto\ (\Pe,-F,Z_1).
\end{equation}
Under \eqref{eq:subst} the objective $\tr(Q\Sig)+\tr(RZ_0)$ becomes
$\tr(W\Pe)+\tr(VZ_1)$, and the constraint blocks
$\mathcal G_{\mathrm c}$ and $\begin{bsmallmatrix}Z_0&L\\L^\top&\Sig
\end{bsmallmatrix}$ become $\mathcal G_{\mathrm e}$ and
$\begin{bsmallmatrix}Z_1&F^\top\\F&\Pe\end{bsmallmatrix}$, respectively. The
system-theoretic hypotheses map correctly:
\[
\text{$(A_\Th,B_\Th)$ stabilizable}\ \leftrightarrow\
\text{$(A_\Th^\top,C_\Th^\top)$ stabilizable}\ \equiv\
\text{$(A_\Th,C_\Th)$ detectable},
\]
and $(A_\Th,Q^{\half})$ detectable $\leftrightarrow$
$(A_\Th^\top,W^{\half})$ detectable $\equiv$ $(A_\Th,W^{\half})$ stabilizable,
which is \emph{automatic} because $W\succ0$ makes $(A_\Th,W^{\half})$
controllable. The role of ``$W\succ0$'' in the control proof (guaranteeing
$\Sig^\star\succ0$ and forcing the closed loop Schur) is played here by
``$Q\mapsto W\succ0$'', and the role of ``$R\succ0$'' by ``$V\succ0$''. Hence:

\begin{corollary}[Estimation SDP]\label{cor:est}
Under Assumption~\ref{as:sys}, all conclusions of
Theorem~\ref{thm:main}(i)--(iv), Lemma~\ref{lem:slater-cont},
Proposition~\ref{prop:primalopt}, Proposition~\ref{prop:c1value}, and
Theorem~\ref{thm:invariant} hold for the estimation SDP \eqref{eq:Jest} under the
substitution \eqref{eq:subst}. In particular \eqref{eq:Jest} has a Slater point,
a unique primal optimizer $\Pe{}^\star\succ0$ equal to the stabilizing solution
of the filter DARE
\[
\Pe=A_\Th^\top\Pe A_\Th
-A_\Th^\top\Pe C_\Th^\top(V+C_\Th\Pe C_\Th^\top)^{-1}C_\Th\Pe A_\Th+W,
\]
a $C^1$ value, and a dual-invariant envelope gradient; consequently the
estimation-side design gradients of Corollary~1 \textup{(}the parallel of
Eqs.~(12)--(13) in $C_\Th,W,V$, and Eq.~(17)\textup{)} read from the returned LMI
duals are exact.
\end{corollary}

\begin{proof}
Every step of Sections~\ref{sec:convex}--\ref{sec:a3} is invariant under the
linear relabeling \eqref{eq:subst}; the hypotheses required there ---
stabilizability of the ``$(A,B)$'' pair, detectability of the ``$(A,Q^{\half})$''
pair, and positive definiteness of the ``$W$'' and ``$R$'' data --- are exactly
the detectability of $(A_\Th,C_\Th)$, the (automatic) stabilizability of
$(A_\Th,W^{\half})$, and $Q\mapsto W\succ0$, $R\mapsto V\succ0$ supplied by
Assumption~\ref{as:sys}. The claims transfer verbatim.
\end{proof}

%=======================================================================
\section{Discussion}\label{sec:discuss}
%=======================================================================

Theorem~\ref{thm:main} converts Corollary~1 into a statement under textbook
hypotheses. The logical structure is worth isolating:

\begin{itemize}
\item \textbf{Slater (A1) $\Leftarrow$ stabilizability/detectability + $W,V\succ0$.}
The Slater point is not generic: it is built from a stabilizing feedback (control
side) or output injection (estimation side), i.e.\ from exactly the
controllability/observability data. Positive-definite process/measurement noise
provides the strict margin.
\item \textbf{Uniqueness (A2) $\Leftarrow$ the DARE.} The covariance-form SDP is
an exact convex lift of the LQR/Kalman problem; its unique optimizer is the
Riccati solution. Detectability of $(A_\Th,Q^{\half})$ (automatic if $Q\succ0$)
and $(A_\Th,C_\Th)$ pins the optimizer.
\item \textbf{Gradient correctness needs \emph{less} than (A3).} Because the
value $\tr(PW)$ is a smooth function of the Riccati solution
(Lemma~\ref{lem:daresmooth}), the value is $C^1$ regardless of dual
non-uniqueness, and the envelope right-hand side is dual-invariant
(Theorem~\ref{thm:invariant}). This is why the finite-difference check in the
paper always matched the dual-read gradient, and why the method is robust even at
solutions where strict complementarity is borderline.
\item \textbf{(A3) is the exact upgrade to dual uniqueness.} Strict
complementarity, which reduces here to $\rank S^\star=n$ at the (always rank-$n$)
active Lyapunov LMI, is necessary and sufficient in the AHO sense for the dual to
be a singleton, recovering Assumption~1 as written.
\end{itemize}

In short: \emph{controllability and observability (stabilizability/detectability),
with positive-definite noise and control weights, imply Assumption~1(A1)--(A2)
and the differentiability and exactness of the design gradient; strict
complementarity, verifiable a~posteriori, supplies the remaining (A3).} Theorem~1
therefore applies to the LQG co-design SDPs under standard system-theoretic
hypotheses, and Corollary~1 holds unconditionally under (S1)--(S3).

%=======================================================================
\begin{thebibliography}{9}
%=======================================================================

\bibitem{AHO}
F.~Alizadeh, J.-P.~A.~Haeberly, and M.~L.~Overton,
\emph{Complementarity and nondegeneracy in semidefinite programming},
Mathematical Programming \textbf{77} (1997), 111--128.

\bibitem{Shapiro97}
A.~Shapiro,
\emph{First and second order analysis of nonlinear semidefinite programs},
Mathematical Programming \textbf{77} (1997), 301--320.

\bibitem{BonnansShapiro}
J.~F.~Bonnans and A.~Shapiro,
\emph{Perturbation Analysis of Optimization Problems},
Springer, New York, 2000.

\bibitem{BoydVandenberghe}
S.~Boyd and L.~Vandenberghe,
\emph{Convex Optimization},
Cambridge University Press, 2004.

\bibitem{BoydLMI}
S.~Boyd, L.~El~Ghaoui, E.~Feron, and V.~Balakrishnan,
\emph{Linear Matrix Inequalities in System and Control Theory},
SIAM, Philadelphia, 1994.

\bibitem{AndersonMoore}
B.~D.~O.~Anderson and J.~B.~Moore,
\emph{Optimal Control: Linear Quadratic Methods},
Prentice Hall, Englewood Cliffs, NJ, 1990.

\bibitem{KwakernaakSivan}
H.~Kwakernaak and R.~Sivan,
\emph{Linear Optimal Control Systems},
Wiley-Interscience, New York, 1972.

\bibitem{LancasterRodman}
P.~Lancaster and L.~Rodman,
\emph{Algebraic Riccati Equations},
Oxford University Press, 1995.

\end{thebibliography}

\end{document}
